% !TEX root = ../main.tex

この部では、圏論と Lean の考え方を、日常的な設計、リファクタリング、チーム開発に接続します。
ここまでの章では、仕様証明やマイグレーションのように比較的はっきりした対象を扱ってきました。
しかし実務では、さらに広い範囲で「変更しても意味が保たれるか」「部品を安全に合成できるか」「設計上の約束が守られているか」を判断する必要があります。

圏論の強みは、合成を中心に物事を見る点にあります。
部品が正しいだけでは十分ではありません。
部品同士をつないだときに意味が保たれるか、つなぎ方を変えても結果が変わらないか、抽象化の境界を越えたときに必要な性質が保存されるかを考える必要があります。
これは、関数パイプライン、ETL、middleware、service adapter、設定マージ、ログ集約、ドメイン層とインフラ層の変換など、多くの設計問題に現れます。

Lean は、このような設計判断のすべてを自動で保証する道具ではありません。
しかし、重要な一部を小さなモデルとして切り出し、「この変換はこの性質を保つ」「このリファクタリングは結果を変えない」「この adapter は業務ロジックと可換である」といった主張を検査できます。
その証明は、設計文書やレビューコメントよりも厳密で、通常のテストよりも一般的な証拠になります。

この部では、合成可能な設計、リファクタリングの等式化、仕様、実装、テスト、証明の接続、チーム開発での運用を扱います。
特に重要なのは、Lean で証明する範囲を現実的に絞ることです。
すべてを形式化するのではなく、壊れると困る核となる性質を選び、CI やレビューの中で使える形にします。

この部を読み終えるころには、読者は圏論と Lean を「勉強した概念」としてではなく、設計レビューや変更管理に使える判断軸として捉えられるようになるでしょう。

\section*{この部で扱うこと}

\begin{itemize}
\item 合成可能な設計の考え方
\item 入出力型と処理連結の設計
\item リファクタリングを等式として扱う方法
\item effectful code の整理とモナド則
\item 仕様、実装、テスト、証明の役割分担
\item Lean ファイルを設計文書やCIに組み込む方法
\item 証明済みと言える範囲の明示
\end{itemize}
